% !TEX root = ../main.tex
\chapter{Viscoelastic Fluid Mechanics}
\label{chap:fluidodynamic_background}

\providecommand{\mathbb}[1]{\mathbf{#1}}

In this chapter, the continuum mechanics framework governing the flow of incompressible viscoelastic fluids is rigorously established. The primary objective is to derive the complete set of differential operators and constitutive equations that define the physical domain and formulate the residuals for the Physics-Informed Neural Network (PINN) architecture. We begin from the universal conservation principles of mass and linear momentum in three dimensions, and subsequently introduce the physical and geometric assumptions characterizing micro-scale polymer flows in enclosed geometries. The Oldroyd-B constitutive model is then formulated in detail, emphasizing the requirement of material frame-indifference via the Upper-Convected Time Derivative (UCTD) and discussing the rheological implications of extensional kinematics. A dedicated non-dimensionalization procedure, based on viscous stress scaling, is presented to ensure well-conditioned gradients within numerical optimization. Finally, we formalize the four-roll mill benchmark, describing its hyperbolic stagnation-point kinematics, boundary conditions, and the mathematical treatment of the pressure gauge freedom.

% ============================================================
% 2.1 GENERAL CONSERVATION LAWS
% ============================================================
\section{General Conservation Laws}
\label{sec:general_conservation_laws}

From a continuum mechanics perspective, the spatial and temporal motion of a fluid continuum occupying an open domain $\Omega \subset \mathbb{R}^3$ is governed by the universal conservation laws of mass and linear momentum \cite{bird1987dynamics}.

\subsection{Principle of Mass Conservation}
\label{subsec:general_mass}
The differential statement of mass conservation for a fluid with local density $\rho(\boldsymbol{x}, t)$ and velocity vector $\boldsymbol{u}(\boldsymbol{x}, t)$ is expressed through the continuity equation:
\begin{equation}
    \frac{\partial \rho}{\partial t} + \nabla \cdot (\rho \boldsymbol{u}) = 0 .
    \label{eq:general_mass}
\end{equation}
Equation~\eqref{eq:general_mass} dictates that the instantaneous rate of mass accumulation within any arbitrary fixed control volume must be strictly balanced by the net convective mass flux through its bounding control surface.

\subsection{Conservation of Linear Momentum}
\label{subsec:general_momentum}
Applying Newton's second law to an infinitesimal fluid parcel yields the Cauchy momentum equation in Eulerian coordinates:
\begin{equation}
    \rho \left( \frac{\partial \boldsymbol{u}}{\partial t} + \boldsymbol{u} \cdot \nabla \boldsymbol{u} \right) = \nabla \cdot \boldsymbol{\sigma} + \rho \boldsymbol{g} ,
    \label{eq:general_momentum}
\end{equation}
where $\boldsymbol{g}$ denotes the acceleration due to external body forces (gravity), and $\boldsymbol{\sigma}$ is the symmetric second-order Cauchy stress tensor. The left-hand side of Equation~\eqref{eq:general_momentum} represents the material (substantial) acceleration $\mathrm{D}\boldsymbol{u}/\mathrm{D}t$, which combines the local unsteady rate of change $\partial \boldsymbol{u}/\partial t$ with the non-linear convective acceleration term $\boldsymbol{u} \cdot \nabla \boldsymbol{u}$ stemming from spatial velocity gradients.

To distinguish between isotropic pressure forces and internal shear/normal deviatoric stresses, the Cauchy stress tensor is decomposed as:
\begin{equation}
    \boldsymbol{\sigma} = -p \boldsymbol{I} + \boldsymbol{T} ,
    \label{eq:cauchy_stress_decomp}
\end{equation}
where $p$ is the isotropic thermodynamic pressure, $\boldsymbol{I}$ is the identity tensor, and $\boldsymbol{T}$ is the extra-stress (or deviatoric stress) tensor. The extra-stress tensor encapsulates all viscous dissipation and elastic energy storage mechanisms originating from the deformation of the fluid microstructure. Substituting Equation~\eqref{eq:cauchy_stress_decomp} into Equation~\eqref{eq:general_momentum} gives:
\begin{equation}
    \rho \left( \frac{\partial \boldsymbol{u}}{\partial t} + \boldsymbol{u} \cdot \nabla \boldsymbol{u} \right) = -\nabla p + \nabla \cdot \boldsymbol{T} + \rho \boldsymbol{g} .
    \label{eq:general_momentum_expanded}
\end{equation}

% ============================================================
% 2.2 PHYSICAL ASSUMPTIONS AND GOVERNING SYSTEM REDUCTION
% ============================================================
\section{Physical Assumptions and Governing System Reduction}
\label{sec:physical_assumptions}

The full three-dimensional, unsteady equations of motion are intractable without specific constitutive closures and scale-appropriate reductions. For the analysis of polymer solutions within the four-roll mill geometry, we introduce the following physical assumptions:

\begin{enumerate}
    \item \textbf{Single-phase continuum}: The polymer solution is treated as an isotropic, single-phase, homogeneous continuum. Individual polymer macromolecules and solvent molecules are not resolved discretely; rather, their collective thermodynamic and rheological effects are captured through effective macroscopic transport properties.

    \item \textbf{Incompressibility and solenoidal velocity field}: Under the operating pressures and low flow velocities characteristic of liquid polymer systems, density variations are negligible ($\rho = \text{const}$). The continuity equation~\eqref{eq:general_mass} reduces to the solenoidal constraint:
    \begin{equation}
        \nabla \cdot \boldsymbol{u} = 0 .
        \label{eq:incompressibility}
    \end{equation}
    A solenoidal velocity field implies that the velocity vector is divergence-free throughout the entire flow domain. Physically, this means that every material parcel preserves its volume identically during deformation (isochoric motion), and the flow possesses neither volumetric mass sources nor sinks.

    \item \textbf{Isothermal conditions}: Viscous dissipation in the present low-shear regime produces negligible internal temperature rises. Thermal uniformity is maintained, decoupling the energy conservation equation from momentum and allowing fluid properties (viscosities, relaxation time) to be treated as spatially invariant constants.

    \item \textbf{Negligible gravitational body forces}: In micro- and milli-fluidic geometries with characteristic length scales on the order of millimeters to centimeters ($L \sim 10^{-2}\,\mathrm{m}$) and viscous working fluids ($\mu \sim 0.1 - 1\,\mathrm{Pa}\cdot\mathrm{s}$), gravitational body forces are completely overshadowed by viscous and viscoelastic stress gradients. This is confirmed by evaluating the ratio of hydrostatic pressure gradients $\rho g$ to characteristic viscous shear forces $\mu_0 U / L^2$:
    \begin{equation}
        \frac{\rho g L^2}{\mu_0 U} \ll 1 .
    \end{equation}
    Consequently, body force contributions are neglected by setting $\boldsymbol{g} = \boldsymbol{0}$.

    \item \textbf{Two-dimensional steady-state kinematics}: The geometry of the four-roll mill features rollers with a high aspect ratio along the axial dimension ($z$-axis). In the transverse mid-plane of the device, edge effects from the top and bottom confinement are negligible, rendering the flow strictly planar:
    \begin{equation}
        \boldsymbol{u} = \begin{pmatrix} u(x, y) \\ v(x, y) \\ 0 \end{pmatrix}, \qquad \frac{\partial (\cdot)}{\partial z} = 0 .
    \end{equation}
    Furthermore, once the boundary roller rotations are brought to constant angular speeds, the fluid reaches an established laminar regime where all local temporal derivatives vanish identically ($\partial(\cdot)/\partial t = \boldsymbol{0}$).

    \item \textbf{Deviatoric stress decomposition}: For dilute and semi-dilute polymer solutions, the total extra-stress tensor $\boldsymbol{T}$ is partitioned into a purely viscous Newtonian solvent contribution and an elastic polymeric contribution $\boldsymbol{\tau}$ \cite{bird1987dynamics, macosko1994rheology}:
    \begin{equation}
        \boldsymbol{T} = 2\mu_s \boldsymbol{D} + \boldsymbol{\tau} ,
        \label{eq:stress_decomposition}
    \end{equation}
    where $\mu_s$ represents the dynamic viscosity of the Newtonian solvent, and $\boldsymbol{D}$ denotes the symmetric rate-of-strain tensor:
    \begin{equation}
        \boldsymbol{D} = \frac{1}{2} \left[ \nabla \boldsymbol{u} + (\nabla \boldsymbol{u})^T \right] .
        \label{eq:strain_rate_tensor}
    \end{equation}
    This two-fluid stress split reflects the fundamental rheological physics of polymer solutions: the Newtonian solvent responds instantaneously to shear through viscous dissipation ($2\mu_s \boldsymbol{D}$), whereas the flexible polymer coils experience flow-induced stretching and orientation, giving rise to an elastic restoring stress $\boldsymbol{\tau}$ governed by macromolecular relaxation.
\end{enumerate}

Substituting the stress decomposition~\eqref{eq:stress_decomposition} and the physical simplifications into Equation~\eqref{eq:general_momentum_expanded} yields the reduced steady two-dimensional momentum equation:
\begin{equation}
    \rho (\boldsymbol{u} \cdot \nabla \boldsymbol{u}) = -\nabla p + \mu_s \nabla^2 \boldsymbol{u} + \nabla \cdot \boldsymbol{\tau} .
    \label{eq:momentum_reduced}
\end{equation}
Here, the Laplacian term $\mu_s \nabla^2 \boldsymbol{u}$ arises directly from $\nabla \cdot (2\mu_s \boldsymbol{D}) = \mu_s \nabla^2 \boldsymbol{u}$ by virtue of incompressibility ($\nabla \cdot \boldsymbol{u} = 0$). Equation~\eqref{eq:momentum_reduced} represents the Navier--Stokes momentum balance coupled directly to the divergence of the polymeric extra-stress tensor, $\nabla \cdot \boldsymbol{\tau}$.

% ============================================================
% 2.3 VISCOELASTIC CONSTITUTIVE MODELING: THE OLDROYD-B FLUID
% ============================================================
\section{Viscoelastic Constitutive Modeling: The Oldroyd-B Fluid}
\label{sec:oldroyd_b}

The momentum equation~\eqref{eq:momentum_reduced} contains six unknown scalar fields in 2D ($u, v, p, \tau_{xx}, \tau_{xy}, \tau_{yy}$), but provides only three equations (continuity and two momentum components). To mathematically close the boundary value problem, a constitutive equation governing the evolution of the polymeric extra-stress tensor $\boldsymbol{\tau}$ is required.

\subsection{Microscopic Formulation and Dimensional Constitutive Law}
\label{subsec:oldroyd_formulation}
In this work, we adopt the classical Oldroyd-B constitutive model \cite{oldroyd1950formulation}. From a kinetic theory perspective, the Oldroyd-B model represents polymer molecules as dilute, non-interacting Hookean dumbbells immersed in a Newtonian continuous medium. Each dumbbell comprises two Brownian beads connected by a linear, infinitely extensible entropic spring \cite{bird1987dynamics}. The beads experience viscous drag forces from the solvent, while the internal spring exerts an entropic restoring force proportional to the end-to-end vector of the dumbbell.

Under these microstructural assumptions, the polymeric stress tensor $\boldsymbol{\tau}$ is governed by the differential constitutive equation:
\begin{equation}
    \lambda \overset{\nabla}{\boldsymbol{\tau}} + \boldsymbol{\tau} = 2\mu_p \boldsymbol{D} ,
    \label{eq:oldroyd_b_dimensional}
\end{equation}
where $\lambda$ is the characteristic polymer relaxation time, $\mu_p$ is the polymeric viscosity contribution, and $\overset{\nabla}{\boldsymbol{\tau}}$ denotes the objective time derivative. The total zero-shear viscosity of the solution is given by the sum of both phases:
\begin{equation}
    \mu_0 = \mu_s + \mu_p .
    \label{eq:zero_shear_viscosity}
\end{equation}

\subsection{Material Frame-Indifference and the Upper-Convected Time Derivative}
\label{subsec:uctd}
In non-Newtonian fluid mechanics, constitutive equations cannot employ the standard partial time derivative $\partial \boldsymbol{\tau}/\partial t$ nor the substantial derivative $\mathrm{D}\boldsymbol{\tau}/\mathrm{D}t$, as these operators are not invariant under time-dependent Euclidean transformations of the observer frame (principle of material frame-indifference or objectivity) \cite{bird1987dynamics, owens2002computational}. 

To account for the continuous translation, rotation, and stretching of fluid elements containing oriented polymer chains, the rate of change of $\boldsymbol{\tau}$ must be evaluated using the \textit{Upper-Convected Time Derivative} (UCTD), denoted by $\overset{\nabla}{\boldsymbol{\tau}}$:
\begin{equation}
    \overset{\nabla}{\boldsymbol{\tau}} \equiv \frac{\partial \boldsymbol{\tau}}{\partial t} + \boldsymbol{u} \cdot \nabla \boldsymbol{\tau} - (\nabla \boldsymbol{u}) \cdot \boldsymbol{\tau} - \boldsymbol{\tau} \cdot (\nabla \boldsymbol{u})^T .
    \label{eq:uctd_definition}
\end{equation}
In Equation~\eqref{eq:uctd_definition}, each term embodies a clear physical mechanism:
\begin{itemize}
    \item $\partial \boldsymbol{\tau}/\partial t$ is the local rate of stress change at a fixed spatial coordinate;
    \item $\boldsymbol{u} \cdot \nabla \boldsymbol{\tau}$ represents the convective advection of stress carried along the flow streamlines;
    \item $-(\nabla \boldsymbol{u}) \cdot \boldsymbol{\tau} - \boldsymbol{\tau} \cdot (\nabla \boldsymbol{u})^T$ are the non-linear tensorial stretching and rotation terms. They quantify how local velocity gradients deform and rotate the macromolecular dumbbell conformations, aligning the polymer stress with the principal axes of extension.
\end{itemize}
For the steady-state regime investigated in this research ($\partial \boldsymbol{\tau}/\partial t = \boldsymbol{0}$), the UCTD reduces strictly to the convective and deformation-induced stretching contributions.

\subsection{Rheological Behavior and the High-Weissenberg Challenge}
\label{subsec:high_weissenberg}
The Oldroyd-B model is widely recognized as the fundamental benchmark for elastic Boger fluids—polymer solutions formulated to exhibit significant elasticity while maintaining an essentially constant shear viscosity over a broad range of shear rates ($\eta(\dot{\gamma}) = \mu_0$). 

However, in strong extensional kinematics (such as those generated along the axes of a four-roll mill), the Oldroyd-B model reveals an intrinsic mathematical singularity. Because Hookean dumbbells possess an infinitely extensible spring law, the steady extensional viscosity $\eta_E$ predicted in a planar extensional flow field ($u = \dot{\epsilon} x, v = -\dot{\epsilon} y$) takes the analytical form:
\begin{equation}
    \eta_E(\dot{\epsilon}) = 4\mu_s + \frac{2\mu_p}{(1 - 2\lambda\dot{\epsilon})(1 + 2\lambda\dot{\epsilon})} .
    \label{eq:extensional_viscosity}
\end{equation}
As the dimensionless extensional rate approaches the critical threshold $2\lambda\dot{\epsilon} \to 1$ (corresponding to an extensional Weissenberg number $Wi \to 0.5$), the polymeric stress diverges asymptotically to infinity ($\eta_E \to \infty$). 

In numerical computations and machine learning architectures alike, this steep stress amplification produces extremely sharp spatial boundary layers along outflow streams, giving rise to the classical \textit{High-Weissenberg Number Problem} (HWNP). Capturing these localized, steep stress gradients without inducing numerical instabilities or gradient explosion is one of the primary motivations for adopting physics-informed neural network approaches.

% ============================================================
% 2.4 NONDIMENSIONALIZATION AND FINAL MODEL EQUATIONS
% ============================================================
\section{Nondimensionalization and Governing Equations}
\label{sec:nondimensionalization}

Nondimensionalization scales variables to order unity, improving numerical conditioning and highlighting the non-dimensional parameters controlling the flow regime \cite{owens2002computational}. In highly viscous and low-Reynolds viscoelastic flows (such as microfluidics and polymer processing in the four-roll mill), choosing a \textit{viscous stress scale} rather than an inertial scale prevents loss of precision and gradient stiffness in the neural network optimization.

\subsection{Reference Scales and Dimensionless Parameters}
Selecting a reference length $H_{\mathrm{ref}}$ (e.g., roller radius), a reference velocity $U_{\mathrm{ref}}$ (roller surface speed), and a constant reference viscosity scale $\eta_0$ (e.g., characteristic total viscosity), we define the dimensionless coordinates, velocities, pressure, and extra-stress tensor:
\begin{equation}
    \boldsymbol{x}^{*} = \frac{\boldsymbol{x}}{H_{\mathrm{ref}}}, \quad
    \boldsymbol{u}^{*} = \frac{\boldsymbol{u}}{U_{\mathrm{ref}}}, \quad
    p^{*} = \frac{p}{\frac{\eta_0 U_{\mathrm{ref}}}{H_{\mathrm{ref}}}}, \quad
    \boldsymbol{\tau}^{*} = \frac{\boldsymbol{\tau}}{\frac{\eta_0 U_{\mathrm{ref}}}{H_{\mathrm{ref}}}} .
    \label{eq:dimless_variables_viscous}
\end{equation}
Scaling both pressure and extra-stress by the viscous stress unit $\eta_0 U_{\mathrm{ref}} / H_{\mathrm{ref}}$ ensures that all stress and pressure terms remain $\mathcal{O}(1)$ in low-Reynolds regimes, eliminating singular perturbation behavior ($1/Re \to \infty$) as $Re \to 0$.

Dropping asterisks for notational simplicity, the dimensionless system is governed by the following parameters:
\begin{align}
    Re_{\mathrm{scale}} &= \frac{\rho U_{\mathrm{ref}} H_{\mathrm{ref}}}{\eta_0} \quad &\text{(Scaling Reynolds number)}, \label{eq:reynolds_scale} \\
    Wi &= \frac{\lambda U_{\mathrm{ref}}}{H_{\mathrm{ref}}} \quad &\text{(Weissenberg number)}, \label{eq:weissenberg} \\
    \mu_s^* &= \frac{\mu_s}{\eta_0} \quad &\text{(Dimensionless solvent viscosity)}, \label{eq:mu_s_star} \\
    \mu_p^* &= \frac{\mu_p}{\eta_0} \quad &\text{(Dimensionless polymeric viscosity)} . \label{eq:mu_p_star}
\end{align}
From these parameters, the physical viscosity ratio $\beta = \mu_s / (\mu_s + \mu_p) = \mu_s^* / (\mu_s^* + \mu_p^*)$ and the physical Reynolds number $Re_{\mathrm{phys}} = \rho U_{\mathrm{ref}} H_{\mathrm{ref}} / (\mu_s + \mu_p) = Re_{\mathrm{scale}} / (\mu_s^* + \mu_p^*)$ can be uniquely determined.

\subsection{Dimensionless Vector System}
The non-dimensional governing equations in steady 2D flow become:
\begin{align}
    \nabla \cdot \boldsymbol{u} &= 0 , \label{eq:dimless_mass} \\
    Re_{\mathrm{scale}} (\boldsymbol{u} \cdot \nabla \boldsymbol{u}) + \nabla p - \mu_s^* \nabla^2 \boldsymbol{u} - \nabla \cdot \boldsymbol{\tau} &= \boldsymbol{0} , \label{eq:dimless_momentum} \\
    Wi \left( \boldsymbol{u} \cdot \nabla \boldsymbol{\tau} - (\nabla \boldsymbol{u})^T \cdot \boldsymbol{\tau} - \boldsymbol{\tau} \cdot \nabla \boldsymbol{u} \right) + \boldsymbol{\tau} - 2\mu_p^* \boldsymbol{D} &= \boldsymbol{0} . \label{eq:dimless_oldroyd}
\end{align}
In this formulation, the constitutive equation~\eqref{eq:dimless_oldroyd} depends strictly on the rheological parameters ($Wi, \mu_p^*$) and is decoupled from inertia ($Re_{\mathrm{scale}}$) and solvent viscosity ($\mu_s^*$), which is crucial for staged parameter identifiability in inverse PINN solvers.

\subsection{Explicit 2D Cartesian PDE Residuals}
Expanding Equations~\eqref{eq:dimless_momentum} and \eqref{eq:dimless_oldroyd} in 2D Cartesian coordinates $(x, y)$ with $\boldsymbol{u} = (u, v)^T$ and $\boldsymbol{\tau} = \begin{pmatrix} \tau_{xx} & \tau_{xy} \\ \tau_{xy} & \tau_{yy} \end{pmatrix}$ yields five scalar PDE residuals. Symmetry of the stress tensor ($\tau_{xy} = \tau_{yx}$) reduces the system to two momentum residuals ($f_u, f_v$) and three constitutive residuals ($f_{\tau_{xx}}, f_{\tau_{yy}}, f_{\tau_{xy}}$):

\begin{align}
    f_u &= Re_{\mathrm{scale}} \left( u \frac{\partial u}{\partial x} + v \frac{\partial u}{\partial y} \right) + \frac{\partial p}{\partial x} - \mu_s^* \left( \frac{\partial^2 u}{\partial x^2} + \frac{\partial^2 u}{\partial y^2} \right) - \left( \frac{\partial \tau_{xx}}{\partial x} + \frac{\partial \tau_{xy}}{\partial y} \right) = 0 , \label{eq:res_fu} \\[1.5ex]
    f_v &= Re_{\mathrm{scale}} \left( u \frac{\partial v}{\partial x} + v \frac{\partial v}{\partial y} \right) + \frac{\partial p}{\partial y} - \mu_s^* \left( \frac{\partial^2 v}{\partial x^2} + \frac{\partial^2 v}{\partial y^2} \right) - \left( \frac{\partial \tau_{xy}}{\partial x} + \frac{\partial \tau_{yy}}{\partial y} \right) = 0 , \label{eq:res_fv}
\end{align}

\begin{equation}
    \begin{aligned}
    f_{\tau_{xx}} = \tau_{xx} + Wi \left( u \frac{\partial \tau_{xx}}{\partial x} + v \frac{\partial \tau_{xx}}{\partial y} - 2 \frac{\partial u}{\partial x} \tau_{xx} - 2 \frac{\partial u}{\partial y} \tau_{xy} \right) - 2\mu_p^* \frac{\partial u}{\partial x} = 0 ,
    \end{aligned}
    \label{eq:res_ftxx}
\end{equation}

\begin{equation}
    \begin{aligned}
    f_{\tau_{yy}} = \tau_{yy} + Wi \left( u \frac{\partial \tau_{yy}}{\partial x} + v \frac{\partial \tau_{yy}}{\partial y} - 2 \frac{\partial v}{\partial x} \tau_{xy} - 2 \frac{\partial v}{\partial y} \tau_{yy} \right) - 2\mu_p^* \frac{\partial v}{\partial y} = 0 ,
    \end{aligned}
    \label{eq:res_ftyy}
\end{equation}

\begin{equation}
    \begin{aligned}
    f_{\tau_{xy}} = \tau_{xy} + Wi \left( u \frac{\partial \tau_{xy}}{\partial x} + v \frac{\partial \tau_{xy}}{\partial y} - \frac{\partial u}{\partial x} \tau_{xy} - \frac{\partial u}{\partial y} \tau_{yy} - \frac{\partial v}{\partial x} \tau_{xx} - \frac{\partial v}{\partial y} \tau_{xy} \right) \\
    - \mu_p^* \left( \frac{\partial u}{\partial y} + \frac{\partial v}{\partial x} \right) = 0 .
    \end{aligned}
    \label{eq:res_ftxy}
\end{equation}

% ============================================================
% 2.5 THE FOUR-ROLL MILL BENCHMARK AND BOUNDARY CONDITIONS
% ============================================================
\section{Four-Roll Mill Boundary Conditions}
\label{sec:four_roll_mill_bcs}

The four-roll mill is a canonical experimental and computational benchmark designed to generate controlled two-dimensional planar extensional flow with a central stagnation point \cite{taylor1934formation, bentley1986experimental}. Originally devised by Taylor \cite{taylor1934formation} to study droplet deformation and breakup under well-defined extensional kinematics, the apparatus has since become a standard test case for validating viscoelastic flow solvers. As illustrated in Figure~\ref{fig:4roll_geometry}, the setup consists of four cylindrical rollers placed symmetrically inside a square domain.

\begin{figure}[htbp]
    \centering
    \IfFileExists{media/4roll_geometry.pdf}{%
        \includegraphics[width=0.75\linewidth]{media/4roll_geometry.pdf}%
    }{%
        \framebox[0.75\linewidth]{\vbox to 5cm{\vfill\centering [Figure Placeholder: media/4roll\_geometry.pdf]\vfill}}%
    }
    \caption{Schematic representation of the four-roll mill geometry, indicating roller counter-rotation speeds $\boldsymbol{\Omega}_k$, the central hyperbolic stagnation point $(0,0)$, and inflow/outflow axes.}
    \label{fig:4roll_geometry}
\end{figure}

\subsection{Geometry and Roller Rotation Dynamics}
The domain consists of a square cavity containing four identical cylindrical rollers of radius $R$ positioned symmetrically. The rollers rotate at equal angular speeds $\Omega$. Adjacent rollers rotate in opposite directions:
\begin{itemize}
    \item Top-Left roller (TL): rotates clockwise ($\boldsymbol{\Omega}_{\mathrm{TL}} = -\Omega \boldsymbol{e}_z$).
    \item Top-Right roller (TR): rotates counter-clockwise ($\boldsymbol{\Omega}_{\mathrm{TR}} = +\Omega \boldsymbol{e}_z$).
    \item Bottom-Left roller (BL): rotates counter-clockwise ($\boldsymbol{\Omega}_{\mathrm{BL}} = +\Omega \boldsymbol{e}_z$).
    \item Bottom-Right roller (BR): rotates clockwise ($\boldsymbol{\Omega}_{\mathrm{BR}} = -\Omega \boldsymbol{e}_z$).
\end{itemize}

This counter-rotating configuration draws fluid inward along the vertical axis (compression axis) and expels it outward along the horizontal axis (extension axis), creating a \textit{hyperbolic} (saddle-type) stagnation point at the origin $(0, 0)$. At this point, the velocity vanishes ($\boldsymbol{u} = \boldsymbol{0}$) while the velocity gradient tensor $\nabla \boldsymbol{u}$ possesses real eigenvalues of equal magnitude and opposite signs, corresponding to simultaneous extension and compression along orthogonal axes. Near the origin, the velocity field approaches pure planar extension:
\begin{equation}
    u(x, y) = \dot{\epsilon}\, x , \qquad v(x, y) = -\dot{\epsilon}\, y ,
    \label{eq:planar_extensional_kinematics}
\end{equation}
where $\dot{\epsilon} \propto \Omega$ is the extensional strain rate. This kinematic configuration makes the four-roll mill particularly relevant for viscoelastic flows, as the persistent extensional deformation near the stagnation point can stretch polymer chains far from their equilibrium configuration, amplifying elastic stress effects.

\subsection{Boundary Conditions for Velocity and Stress}
To complete the boundary value problem in CFD, appropriate boundary conditions are imposed:

\begin{enumerate}
    \item \textbf{No-Slip Velocity BCs on Roller Walls ($\Gamma_{\mathrm{rolls}}$)}:
    On the surface of each roller $k \in \{\mathrm{TL}, \mathrm{TR}, \mathrm{BL}, \mathrm{BR}\}$, the fluid velocity matches the tangential roller wall speed:
    \begin{equation}
        \boldsymbol{u}|_{\Gamma_k} = \boldsymbol{\Omega}_k \times (\boldsymbol{x} - \boldsymbol{x}_k) ,
        \label{eq:bc_velocity_rolls}
    \end{equation}
    where $\boldsymbol{x}_k$ is the center coordinate of roller $k$.

    \item \textbf{Outer Boundary Conditions ($\Gamma_{\mathrm{outer}}$)}:
    On the outer bounding walls of the cavity, no-slip velocity conditions ($\boldsymbol{u} = \boldsymbol{0}$) or prescribed inflow/outflow velocity profiles corresponding to the extensional field~\eqref{eq:planar_extensional_kinematics} are specified.

    \item \textbf{Polymeric Extra-Stress Wall BCs}:
    For viscoelastic fluids, stress boundary conditions on solid walls must be consistent with the local kinematics. On the rotating rollers, the stress tensor $\boldsymbol{\tau}$ satisfies either the local steady-state constitutive relation evaluated at the wall velocity gradient $\nabla \boldsymbol{u}|_{\Gamma_{\mathrm{rolls}}}$ or is constrained via Dirichlet conditions extracted from COMSOL reference simulations.
\end{enumerate}

\subsection{Pressure Pinning and Gauge Fix in Enclosed Domains}
In incompressible fluid mechanics, the pressure $p$ appears in the momentum equation solely through its spatial gradient $\nabla p$. As a consequence:
\begin{itemize}
    \item If a velocity boundary condition (Dirichlet no-slip) is specified along the entire closed boundary $\partial \Omega = \Gamma_{\mathrm{rolls}} \cup \Gamma_{\mathrm{outer}}$, the pressure field is determined only up to an arbitrary additive constant $C$:
    \begin{equation}
        p(\boldsymbol{x}) \to p(\boldsymbol{x}) + C .
    \end{equation}
    \item Without a pressure reference point, the continuous Neumann problem for pressure lacks a unique solution, causing numerical solvers (both traditional FEM/FVM and PINN optimization) to suffer from severe pressure drift.
\end{itemize}

To ensure mathematical well-posedness and unique pressure recovery, a single-point pressure pin (gauge condition) is enforced at a reference location $\boldsymbol{x}_0$ (typically the stagnation point or a domain corner):
\begin{equation}
    p(\boldsymbol{x}_0) = p_0 = 0 .
    \label{eq:pressure_pin}
\end{equation}
Fixing $p(\boldsymbol{x}_0) = 0$ removes the hydrostatic gauge ambiguity, allowing the numerical solver to uniquely reconstruct the absolute pressure field across the four-roll mill domain. In the PINN framework, this gauge condition is enforced as an additional penalty term in the composite loss function, $L_{\mathrm{pin}} = |\hat{p}(\boldsymbol{x}_0) - p_0|^2$, which anchors the predicted pressure to the reference value during optimization.
